\section{Monitoreo del entrenamiento y sistema de registro}

\subsection{Forma estructurada de escribir el registro}

Para poder reproducir todo el proceso de entrenamiento es fundamental contar con un logging schema detallado. Los siguientes campos son el requisito mínimo:

\begin{itemize}
    \item `step`, `episode`, `wall_time`: para ubicar el slice en que ocurrió el problema
    \item `avg_return`, `std_return`, `kl_divergence`, `entropy`: para medir la tendencia futura de la policy
    \item `buffer_size`, `td_error_mean`, `td_error_std`: el estado interno del Replay buffer
    \item `learning_rate`, `grad_norm`, `value_loss`: para diagnosticar el optimizador y el value head
\end{itemize}

\begin{importantbox}{Granularidad del registro y alarmas}
El logging de alta frecuencia (cada 1k step) combinado con threshold-based alarms (KL > 0.06, return variance > 0.2) puede interrumpir el entrenamiento antes de que ocurra una divergence, lo cual es particularmente importante en entornos Sim2Real para evitar un daño físico.
\end{importantbox}

\subsection{Panel de métricas}

El dashboard de visualización en tiempo real generalmente incluye al menos:

\begin{enumerate}
    \item PPO ratio distribution histogram
    \item SAC entropy \& Q-value curves
    \item Batch-wise advantage histogram to detect saturation
\end{enumerate}

\subsection{Resumen de la sección}

Un logging estructurado junto con un dashboard en tiempo real les permite a los investigadores detectar a tiempo failure modes como divergence, explosion y overfitting, lo cual permite reproducir más rápidamente el desempeño central.

